\documentclass[article]{jss}

\usepackage{amsmath}
\usepackage{amssymb}
\usepackage{amsfonts}
\usepackage{graphicx}
\usepackage{url}
\usepackage{hyperref}
\usepackage{booktabs}
\usepackage{color}
\usepackage{natbib}

%% Fallbacks so the paper also builds against a minimal jss.cls.
\providecommand{\proglang}[1]{\textsf{#1}}
\providecommand{\pkg}[1]{\textbf{#1}}
\providecommand{\code}[1]{\texttt{#1}}

\setlength{\emergencystretch}{3em}

\author{
  Peter Cotton\thanks{\email{peter.cotton@microprediction.com}}\\
  \url{https://skaters.microprediction.org}
}

\title{A Stronger Univariate Benchmark for Foundational Time-Series Model Evaluation}
\Plaintitle{A Stronger Univariate Benchmark for Foundational Time-Series Model Evaluation}
\Shorttitle{A Stronger Univariate Benchmark}

\newcommand{\paperabstract}{
A foundation model forecasts from a fixed window, with weights frozen at
pretraining. An online forecaster carries state and revises it at every
observation. Benchmarks for the first kind contain no example of the second: of
the 132 entries on the \proglang{GIFT-Eval} leaderboard, five are not neural, and
all five refit from scratch on each window. We propose \code{laplace}, a
zero-dependency online distributional forecaster with no pretraining and no
tunable parameters, as the missing reference. Scored one step ahead on
held-out log density over an 18{,}815-series \proglang{FRED} panel, it beats
TimesFM 3.0 on the median series in every stratum by roughly half a nat, and a
post-hoc portfolio that sees only the model's own issued density recovers most
of the gap without retraining. The reference costs a \code{pip install} to include, and a
\proglang{Rust} core with a bit-exact parity suite makes its numbers recompute
on other hardware.
}

\begin{document}

\maketitle

\begin{abstract}
\paperabstract
\end{abstract}

\noindent\textbf{Keywords:} time series, foundation models, benchmarking,
probabilistic forecasting, proper scoring rules, online learning.

%% -------------------------------------------------------------------------
\section{Introduction}\label{sec:intro}

A foundation model forecasts from a fixed window of recent history, with weights
frozen at pretraining. An online forecaster carries state, and revises it with
every observation it sees. Benchmarks for the first kind contain no example of
the second.

Table~\ref{tab:arms} lists the non-neural arms of the studies that rank these
models. The set barely moves between them: naive, seasonal naive, and the three
\pkg{statsforecast} calls, sometimes with the \pkg{SCUM} ensemble of
\citet{petropoulos2020scum}. Each study also carries deep models, but those are
trained per dataset and serve as competitors rather than as references.

\begin{table}[t!]
\centering
\begin{tabular}{llcc}
\toprule
Study & Non-neural arms & Online & Density \\
\midrule
\citet{aksu2024gifteval}    & Naive, SNaive, AutoARIMA, AutoETS, AutoTheta       & no & no \\
\citet{ansari2024chronos}   & the above, plus the \pkg{SCUM} ensemble             & no & no \\
\citet{ansari2025chronos2}  & the above; task-specific deep models dropped        & no & no \\
\citet{das2024timesfm}      & ETS, ARIMA, \pkg{CatBoost}                          & no & no \\
\citet{auer2025tirex}       & AutoARIMA, AutoTheta, SNaive                        & no & no \\
\citet{woo2024moirai}       & AutoARIMA, SNaive                                   & no & no \\
\citet{shchur2025fevbench}  & Naive, SNaive, Drift, the three above, ensemble     & no & no \\
\bottomrule
\end{tabular}
\caption{Non-neural arms in the studies that rank time-series foundation
models. \emph{Online} means the method carries state between forecasts rather
than refitting. \emph{Density} means it issues a predictive density rather than
a point or a quantile vector. No arm in any study is either.}
\label{tab:arms}
\end{table}

None of these arms is online. The \pkg{statsforecast} predict path is handed an
array, fits, forecasts, and keeps nothing, so every window starts over. In
\citet{ansari2024chronos} and in the Monash tables of \citet{das2024timesfm}
there is a single forecast origin per series, so the question does not arise.

None of them issues a density either. Their quantiles come from a Gaussian
residual-variance interval, which the harness reads off level by level, so the
predictive is symmetric by construction with no skew and no tail. The ensemble
is worse: its quantiles are a per-level median across members, which is not a
distribution at all.

\pkg{NPTS} would fill the gap. It is training-free, fully distributional, and
ships inside \pkg{GluonTS}, which these harnesses already import. It appears in
none of the studies in Table~\ref{tab:arms}.

This paper supplies such an arm and reports what it does to the comparison. The
contribution is deliberately small: a reference, a protocol, and four tables.

%% -------------------------------------------------------------------------
\section{\code{laplace}: an online distributional forecaster}\label{sec:reference}

\code{laplace} imports nothing, requires no GPU, performs no fitting step, and
exposes no tunable hyperparameters. One parameter-free call returns a predictive
density.

The predictive object is a weighted mixture of Gaussians. That choice buys
three things a reference arm needs: the mixture is closed under affine maps and
under mixing, it admits a closed-form CRPS, and it approximates heavy tails as a
scale mixture. A density can therefore be pushed through an invertible map and
scored against a proper rule \citep{gneiting2007strictly} without sampling.

A \emph{transform} is a stateful causal map that consumes one observation and
emits one transformed value, paired with an inverse that rebuilds predictive
distributions for future observations from the transformed coordinate, given the
state it has accumulated. Differencing, standardisation, online autoregression,
\pkg{GARCH}-style scaling, seasonal differencing and the Yeo--Johnson
coordinates are all transforms. \emph{Conjugation} composes one with an inner
forecaster: the observation is transformed, the inner model issues a density in
the new coordinate, and the inverse carries that density back. Multi-step
uncertainty accumulates through the inverse rather than by simulation.

\code{laplace} is one fixed composition of these parts, with no knob exposed.
Candidate models enter it only as mean forecasters, weighted online by
predictive likelihood, so that a shaped leaf governs the tail rather than the
dispersion of candidate means. A lattice projection places atoms on the exact
values a grid-quoted series revisits, and vanishes identically on continuous
data. Beyond one step the predictive mixes over decimated clocks, since the fine
clock wins short horizons and a coarse one wins long. A conditional tail fit
splices the extreme quantiles, which a bulk fit leaves too thin.

\citet{cotton2026transforms} derives all of this. We restate only what a reader
needs to judge the comparison.

Two properties matter for a reference arm, and neither is accuracy. The first is
cost. Adding \code{laplace} to a leaderboard is a \code{pip install} and one
function call, with no weights to download, no device to allocate, and no
context length to choose. The pure-\proglang{Python} implementation, $280$\,KB of source across $28$
files, is the reference, and a \proglang{Rust} core steps it at roughly $14$ microseconds
per observation, twenty times the \proglang{Python} rate. A reference that is
slow over a large panel does not get included.

The second is reproducibility across machines. The \proglang{Rust} core uses
portable mathematics (\pkg{libm}) and produces identical bits on x86 and ARM,
and a shared parity suite checks roughly $100{,}000$ probe values to $10^{-6}$
per implementation. A number one person publishes therefore recomputes on
another person's hardware, which is the property a denominator needs and a
pretrained checkpoint behind an inference stack does not have.

A denominator sets what a skill score means, and a weak one flatters everything
divided by it. That is the argument for this particular arm.

Sections~\ref{sec:classic} and \ref{sec:results} measure how strong it is.
\code{laplace} takes the per-series log-likelihood against every classical arm
in Table~\ref{tab:arms}, and against TimesFM 3.0 in every stratum tested.
\citet{cotton2026transforms} reports the same ordering against neural and
foundation-model baselines across the non-price \proglang{FRED} universe;
\pkg{GARCH}-$t$ on asset prices is the one arm that escapes, and that split is
reported rather than averaged away.

%% -------------------------------------------------------------------------
\section{Protocol}\label{sec:protocol}

Every arm sees the same series, the same windows and the same scoring rule. The
context is 128 observations, the test window rolls one step ahead, and no arm is
fitted to the evaluation data. The score is held-out log predictive density,
differenced per series so that each series' intrinsic difficulty cancels.

We report the median of that difference and the per-series win, draw and loss
counts from a Diebold--Mariano test with HAC standard errors and a draw band,
following \citet{diebold1995comparing}. We do not report a pooled mean.
Averaging a log density across heterogeneous series estimates nothing, because a
few near-constant series dominate the sum; \citet{koning2005m3} made the
corresponding argument about the M3 competition, and \citet{hewamalage2023evaluation}
prescribe per-series testing.

The context is fixed rather than growing. Several foundation models degrade
beyond the window they were pretrained on, and most of them silently truncate
over-long input rather than raising, so feeding four thousand points to a model
configured for two thousand scores a two-thousand-context model without saying
so. A protocol that lets context grow needs a declared per-model cap, set at the
documented architectural maximum, and recorded per arm.

Recording it matters for a second reason. The \citet{woo2024moirai} benchmark
numbers come from a per-dataset search over context lengths from one thousand to
five thousand, jointly with patch size, which the arms it is compared against do
not receive.

%% -------------------------------------------------------------------------
\section{Against the classic arms}\label{sec:classic}

The arms of Table~\ref{tab:arms} are the denominator every published skill
score divides by, so their strength sets what a skill score means.
Table~\ref{tab:classic} scores \code{laplace} against them on the non-price
\proglang{FRED} universe.

\begin{table}[t!]
\centering
\begin{tabular}{lccccr}
\toprule
 & \multicolumn{2}{c}{LL win} & \multicolumn{2}{c}{CRPS win} & \\
\cmidrule(lr){2-3}\cmidrule(lr){4-5}
Arm & raw & family & raw & family & $N$ \\
\midrule
AutoARIMA (\pkg{statsforecast}) & $87\,\%$ & $87\,\%$ & $58\,\%$ & $63\,\%$ & $1688$ \\
AutoETS (\pkg{statsforecast})   & $95\,\%$ & $94\,\%$ & $77\,\%$ & $78\,\%$ & $1688$ \\
Theta (\proglang{R})            & $92\,\%$ & $95\,\%$ & $74\,\%$ & $79\,\%$ & $1917$ \\
auto.arima (\proglang{R})       & $84\,\%$ & $85\,\%$ & $51\,\%$ & $60\,\%$ & $1917$ \\
\bottomrule
\end{tabular}
\caption{Share of series on which \code{laplace} scores better, over $5{,}397$
continuous non-price \proglang{FRED} series in $1{,}213$ correlated families.
The family figure averages within a family first, so that a cluster of related
series counts once. $N$ is the subset each arm covers under its compute budget.}
\label{tab:classic}
\end{table}

On log-likelihood the margin is wide: the closest arm loses on $84\,\%$ of
series and the widest on $95\,\%$. On CRPS the arms do better, which is the
expected pattern, since CRPS is forgiving of a mis-shaped tail that a log score
punishes.

This is the sense in which the incumbent denominators are soft. A skill score
computed against them credits a model for clearing a bar that a dependency-free
online method clears by a wide margin on the same data.

\section{Against TimesFM 3.0}\label{sec:results}

Table~\ref{tab:raw} scores TimesFM 3.0 \citep{das2024timesfm} zero-shot against
the reference on an 18{,}815-series \proglang{FRED} panel. Negative favours the
reference.

\begin{table}[t!]
\centering
\begin{tabular}{lrccr}
\toprule
Stratum & $N$ & Win / draw / loss & Median $\Delta$LL & Loss rate \\
\midrule
Daily, economic       & $2{,}179$ & $66 / 923 / 1{,}190$      & $-0.854$ & $54.6\,\%$ \\
Daily, price/returns  & $7{,}633$ & $15 / 4{,}593 / 3{,}025$  & $-0.616$ & $39.6\,\%$ \\
Weekly, economic      & $3{,}013$ & $517 / 1{,}481 / 1{,}015$ & $-0.520$ & $33.7\,\%$ \\
Monthly, economic     & $5{,}576$ & $193 / 3{,}592 / 1{,}791$ & $-0.647$ & $32.1\,\%$ \\
M4-hourly, seasonal   & $414$     & $8 / 241 / 165$           & $-0.543$ & $39.9\,\%$ \\
\bottomrule
\end{tabular}
\caption{TimesFM 3.0 against \code{laplace}, one step ahead on held-out log
density, 128-length context. $\Delta$LL is the per-series difference, model
minus reference, reported as a median. The loss rate is the share of series on
which the model loses.}
\label{tab:raw}
\end{table}

The model loses on the majority of series in every stratum, by a median of
roughly half a nat. Half a nat is a large gap in log-likelihood terms.

Table~\ref{tab:wrapped} repeats the comparison with the model wrapped in a
post-hoc portfolio that combines its own issued predictive with the reference.
The wrap requires no retraining and no access to model internals; it sees only
the density the model published.

\begin{table}[t!]
\centering
\begin{tabular}{lrccr}
\toprule
Stratum & $N$ & Win / draw / loss & Median $\Delta$LL & Loss rate \\
\midrule
Daily, economic       & $2{,}179$ & $76 / 1{,}383 / 720$     & $-0.094$ & $33.0\,\%$ \\
Daily, price/returns  & $7{,}633$ & $152 / 7{,}136 / 345$    & $-0.008$ & $4.5\,\%$ \\
Weekly, economic      & $3{,}011$ & $467 / 2{,}045 / 499$    & $-0.049$ & $16.6\,\%$ \\
Monthly, economic     & $5{,}195$ & $350 / 4{,}480 / 365$    & $-0.056$ & $7.0\,\%$ \\
M4-hourly, seasonal   & $414$     & $58 / 352 / 4$           & $-0.007$ & $1.0\,\%$ \\
\bottomrule
\end{tabular}
\caption{The same model wrapped in a never-worse portfolio with the reference.
The loss rate falls by as much as thirty-nine points; the median stays negative
in every stratum.}
\label{tab:wrapped}
\end{table}

Two readings follow. The wrapped model is safe to deploy, since its loss rate
collapses. It is still not competitive, since the median remains negative
everywhere.

The second reading is the more useful one. If a post-hoc combination that sees
only the published density recovers most of the gap, then at this horizon the
pretraining is contributing less than the raw comparison suggests, and what it
does contribute survives a wrapper.

%% -------------------------------------------------------------------------
\section{The reconstruction moves the answer more than the model does}\label{sec:nozzle}

A quantile-only model does not emit a density, and a density must exist before a
log score does. Something must build one from the quantiles, and that something
is a modelling choice independent of the model being scored.

Fix TimesFM 3.0 and fix $3{,}200$ held-out test points, then vary only the
reconstruction across four methods. The mean log density moves from $-0.373$ to
$+2.104$. The mean per-point spread is $3.13$ nats, which exceeds the entire
measured gap between the model and the reference in every stratum of
Table~\ref{tab:raw}. Two methods already in independent use elsewhere differ by
$0.62$ nats.

The choice therefore belongs in the protocol, stated rather than assumed. It is
not peculiar to this study: the \citet{aksu2024gifteval} submission interface
requires a mean and nine deciles, so a model emitting a genuine density must
reduce it before scoring and a model emitting samples must convert.

The conversion has teeth. That interface's interval score asks for the
$2.5$th and $97.5$th percentiles, and a sample forecast returns a nearest order
statistic without interpolation, so at the twenty samples the reference
\pkg{Chronos} notebook uses, those two quantiles are the minimum and the maximum
of twenty draws.

%% -------------------------------------------------------------------------
\section{Horizon, grid, and statefulness}\label{sec:suites}

The results above do not contradict any published leaderboard, because no
published leaderboard occupies this regime. Three settings separate them.

The horizon never reaches one step. The shortest prediction length in
\citet{aksu2024gifteval} is six, and \citet{shchur2025fevbench} use domain
horizons such as thirty and one hundred sixty-eight.

The scoring rule is a nine-point grid. What these suites report as CRPS is a
mean of weighted quantile losses at the deciles, so nothing outside the tenth
and ninetieth percentile enters any published score.
\citet{ansari2025chronos2} call that grid a defect and move to twenty-one
levels including the first and ninety-ninth.

And no baseline carries state, which is the hole of Section~\ref{sec:intro}.
Each setting is defensible for the task those suites set themselves. Together
they define a regime, and an online distributional method cannot compete inside
it or be seen to.

%% -------------------------------------------------------------------------
\section{Scope}\label{sec:scope}

Everything measured here is univariate and one step ahead. A single pretrained
model amortized across many related series can exploit cross-series and panel
structure that no univariate method reaches by construction, and that is where
pretraining should pay for itself most directly. This paper does not make that
comparison, and it is the more decisive one.

The protocol is assembled from standard parts. Proper scoring rules are the norm
in this literature; prequential test-then-train evaluation is textbook; paired
per-series testing was urged on the field by \citet{koning2005m3} and is already
the headline aggregation in \citet{shchur2025fevbench}; normalizing by a cheap
reference is universal. The assembly is what is new.

Leakage is controlled in this literature by dataset name rather than by time.
\citet{ansari2024chronos} state the temporal standard, that evaluation data
should begin after the last pretraining observation, and then decline it as a
minimal risk. \citet{meyer2025leakage} and \citet{oreshkin2026simulation} report
that a corpus used in TimesFM pretraining reappears in its Monash evaluation;
their pretraining table lists Traffic at $862$ series and $15{,}122{,}928$
observations, and the Monash archive file \citep{zenodo-traffic-hourly} holds
$862$ series of exactly $17{,}544$ steps.

A calibration framing would have to engage \citet{adler2026calibrated}, who find
these models better calibrated than automated ARIMA and N-BEATS on standard
datasets. Nothing here contradicts that. Calibration on a decile grid and
sharpness of a full density are different measurements.

Finally, the panel is \proglang{FRED}, and \proglang{FRED} is the universe these
models' pretraining corpora almost entirely omit: the corpus of
\citet{woo2024moirai} is $0.09\,\%$ economic by observation, and the entire
economics domain of \citet{aksu2024gifteval} is the M4 collection. Whether the
gap survives on the energy and transport data these models were built for is an
open question, and \citet{wang2026spectral} suggest it will not.

%% -------------------------------------------------------------------------
\section{Reproduction}\label{sec:repro}

The paper directory holds four files. \code{constants.py} derives every number
in Sections~\ref{sec:results} and \ref{sec:nozzle} from a version-controlled
store. \code{verify\_paper.py} re-derives each one, compares it against the
tables above, and exits non-zero on drift. \code{REPRODUCE.md} gives three
tiers: the figures check in seconds from committed files, the summaries rebuild
from per-step records in minutes, and the records regenerate from raw data in
days.

The data vintage is published as a release asset in the repository
\code{skaters-\allowbreak benchmarking-\allowbreak data}. \proglang{FRED} revises series, so a fresh
pull returns different vintages and will not reproduce these numbers. The archive carries $210{,}574$ public-domain series; commercial index
series are excluded and listed, which costs the price stratum and nothing else.

%% -------------------------------------------------------------------------
\bibliographystyle{plainnat}
\bibliography{benchmark}

\end{document}
